\documentclass[11pt]{amsart}
\usepackage{amssymb,mathtools,xcolor}
\usepackage[colorlinks=true,linkcolor=blue,citecolor=blue,urlcolor=blue]{hyperref}

\newtheorem{theorem}{Theorem}[section]
\newtheorem{proposition}[theorem]{Proposition}
\newtheorem{lemma}[theorem]{Lemma}
\newtheorem{corollary}[theorem]{Corollary}
\theoremstyle{definition}
\newtheorem{definition}[theorem]{Definition}
\theoremstyle{remark}
\newtheorem{remark}[theorem]{Remark}

\newcommand{\Z}{\mathbb{Z}}
\newcommand{\Q}{\mathbb{Q}}
\newcommand{\C}{\mathbb{C}}
\newcommand{\cO}{\mathcal{O}}
\newcommand{\fl}[1]{\left\lfloor #1 \right\rfloor}
\newcommand{\dpow}[2]{#1^{[#2]}}
\newcommand{\todo}[1]{\textcolor{red}{[TODO: #1]}}
\DeclareMathOperator{\ord}{ord}

\title[Valuations of Romik's sequence]{The $p$-adic valuation of Romik's sequence\\ at primes $p\equiv 3 \pmod 4$}
\author{Fengbo Yuan}
\address{Department of Physics, University of Alabama at Birmingham}
\email{yuanf@uab.edu}
\date{Draft of 22 November 2026}

\begin{document}

\begin{abstract}
Romik's sequence $(d(n))_{n\ge 0}=1,1,-1,51,849,-26199,\dots$ records the Taylor
coefficients of the Jacobi theta constant $\theta_3$ at the CM point $\tau=i$. For a prime
$p\equiv 3\pmod 4$ we prove
$v_p(d(n))\ge v_p\bigl((4n)!\bigr)-v_p\bigl((p\lfloor 4n/(p+1)\rfloor)!\bigr)$,
and we show that equality holds if and only if $p\nmid d(r)$, where $r$ is the residue of $n$
modulo $(p+1)/4$. The lower bound implies a conjecture of Krattenthaler and M\"uller:
$d(n)\equiv 0 \pmod{p^e}$ for all $n\ge \lceil ep^2/4\rceil$. The proof identifies $d(n)$ with
Taylor coefficients of $\sigma(z)\sqrt{\wp(z)}$ on the lemniscatic curve and uses the estimates of
Bannai, Kobayashi and Yasuda for Weierstrass functions at supersingular primes. By an elementary argument we also bound the valuations of Romik's auxiliary
sequences $u(n)$ and $v(n)$, which proves parts (1) and (2) of another conjecture of Krattenthaler
and M\"uller. Finally we show that $d(n)^2$ is, up to an explicit factor, the central value
$L(\psi^{8n+1},4n+1)$ for the Hecke character $\psi$ of the curve $y^2=x^3-x$, so that the results
above give a lower bound for the $p$-adic valuations of the normalised central values at the primes
inert in $\Q(i)$ and characterise when it is attained.
\end{abstract}

\maketitle

\section{Introduction}

Let $\theta_3(x)=\sum_{n\in\Z}e^{-\pi n^2x}$ for $\operatorname{Re}x>0$. Romik \cite{Romik}
considered the function
\[
\sigma_3(z)=\frac{1}{\sqrt{1+z}}\,\theta_3\!\left(\frac{1-z}{1+z}\right)\qquad(|z|<1)
\]
and showed that there is a sequence of integers $(d(n))_{n\ge0}$ with
\begin{equation}\label{eq:romik}
\sigma_3(z)=\theta_3(1)\sum_{n=0}^{\infty}\frac{d(n)}{(2n)!}\,\Phi^n z^{2n},
\qquad \Phi=\frac{\Gamma(1/4)^8}{128\,\pi^4}.
\end{equation}
He conjectured that $d(n)$ vanishes modulo $p$ for all large $n$ when $p\equiv3\pmod4$. This was
proved by Scherer \cite{Scherer} with the threshold $n>(p^2-1)/2$. For prime powers,
Kim and Lee \cite{KimLee} (for $p\ge5$) and Krattenthaler and M\"uller \cite{KM} proved
$d(n)\equiv0\pmod{p^e}$ for $e\ge2$ and $n\ge\lceil (e-1)p^2/2\rceil$, and
\cite[Conjecture~45]{KM} predicts that the threshold $\lceil ep^2/4\rceil$ suffices.

Throughout, $p$ is an odd prime and $v_p$ is the $p$-adic valuation with $v_p(p)=1$.

\begin{theorem}\label{thm:main}
Let $p\equiv3\pmod4$ be a prime and let $n\ge0$. Put $m=\fl{4n/(p+1)}$ and $r=n-m(p+1)/4$, so that
$0\le r<(p+1)/4$.
\begin{enumerate}
\item[(i)] We have
\[
v_p\bigl(d(n)\bigr)\;\ge\; v_p\bigl((4n)!\bigr)-v_p\bigl((pm)!\bigr).
\]
\item[(ii)] Let $\delta(k)=2^k d(k)/(4k)!$. Then $p\,\delta\bigl(\tfrac{p+1}{4}\bigr)$ is a $p$-adic unit and
\[
p^m\,m!\,\delta(n)\;\equiv\;\Bigl(p\,\delta\bigl(\tfrac{p+1}{4}\bigr)\Bigr)^{m} \delta(r)\pmod{p\,\Z_{(p)}} .
\]
In particular, equality holds in \textup{(i)} if and only if $p\nmid d(r)$.
\end{enumerate}
\end{theorem}

Since $v_p((pm)!)=m+v_p(m!)$, part (i) says that $d(n)/(4n)!$ has at worst the denominator of the
coefficient of $z^{(p+1)m}$ in $\exp(z^{p+1}/p)$.

\begin{corollary}\label{cor:KM}
Let $p\equiv3\pmod4$. Then $v_p(d(n))\ge\fl{4n/p^2}$ for all $n\ge0$. In particular
$d(n)\equiv0\pmod{p^e}$ for all $n\ge\lceil ep^2/4\rceil$ and all $e\ge1$, which is
\cite[Conjecture~45]{KM}.
\end{corollary}

\begin{corollary}\label{cor:modp}
Let $p\equiv3\pmod4$. Then $d(n)\equiv0\pmod p$ for all $n>(p^2-1)/4$, and
$p\nmid d\bigl((p^2-1)/4\bigr)$.
\end{corollary}

Theorem~\ref{thm:main} is the specialisation of a statement about Weierstrass functions. For the
Weierstrass $\sigma$-function of an elliptic curve with supersingular reduction, the lower bound
$v_p(a_n)\ge-\tau_{n-1}$ of Theorem~\ref{thm:sigma} below is a theorem of Bannai, Kobayashi and
Yasuda \cite[Theorem~5.6]{BKY}, who prove equality for the orders $n=p^k(p+1)+1$. Our proofs use
their bounds for the coefficients of $\wp$ and $\zeta$ \cite[Corollaries~5.3, 5.5]{BKY}. For the
curves $y^2=4x^3-g_2x$ we determine all orders at which equality holds.

\begin{theorem}\label{thm:sigma}
Let $p\equiv3\pmod4$, let $K$ be a finite unramified extension of $\Q_p$ with ring of integers
$\cO$, and let $g_2\in\cO^\times$. Let $\sigma(z)=\sum_{n\ge1}a_nz^n$ be the Weierstrass
$\sigma$-function of $y^2=4x^3-g_2x$. For $n\ge1$ write $n-1=(p+1)m+s$ with $0\le s\le p$. Then
$p\,a_{p+2}\in\cO^\times$ and
\[
p^m\,m!\,a_n\;\equiv\;(p\,a_{p+2})^m\,a_{s+1}\pmod{p\,\cO}.
\]
In particular $v_p(a_n)\ge-m-v_p(m!)$, with equality if and only if $a_{s+1}\in\cO^\times$; equality
holds for every $n\equiv1\pmod{p+1}$.
\end{theorem}

Romik's recursion for $d(n)$ involves two auxiliary sequences of integers $u(n)$ and $v(n)$. In
the normalisation of \cite[\S2]{KM} they are given by
\begin{equation}\label{eq:uv}
\sum_{n\ge0}\frac{u(n)}{(2n+1)!}\,t^{2n+1}=\frac{t\,F_3(4t^2)}{F_1(4t^2)},\qquad
\sum_{n\ge0}\frac{v(n)}{2^n(2n)!}\,t^{2n}=F_1(4t^2)^{1/2},
\end{equation}
where $F_1(x)={}_2F_1(\tfrac14,\tfrac14;\tfrac12;x)$ and
$F_3(x)={}_2F_1(\tfrac34,\tfrac34;\tfrac32;x)$. For these sequences an elementary argument,
independent of the rest of the paper, gives the following.

\begin{theorem}\label{thm:uv}
Let $p\equiv3\pmod4$ be a prime. For all $n\ge0$,
\[
v_p\bigl(u(n)\bigr)\ge v_p\bigl((2n+1)!\bigr)-\fl{\frac{2n+1}{p+1}},\qquad
v_p\bigl(v(n)\bigr)\ge v_p\bigl((2n)!\bigr)-\fl{\frac{2n}{p+1}} .
\]
\end{theorem}

\begin{corollary}\label{cor:KM43}
Let $p\equiv3\pmod4$. Then
\begin{enumerate}
\item[(i)] $u(n)\equiv0\pmod{p^{2e-1}}$ for all $n\ge\lceil (ep^2-1)/2\rceil$ and all $e\ge1$;
\item[(ii)] $u(n)\equiv0\pmod{p^{2e}}$ for all $n\ge(ep^2+(e-2)p)/2$ and all $e\ge2$;
\item[(iii)] $u(n)\equiv0\pmod{p^2}$ for all $n\ge(p^2-1)/2$, and $u(n)\equiv0\pmod p$ for all
$n\ge(p^2-p)/2$.
\end{enumerate}
\end{corollary}

Parts (i) and (ii) are parts (1) and (2) of \cite[Conjecture~43]{KM}. Not all of the corollary is
new: the cases $e\le2$ of (i) are contained in \cite[Theorem~9]{KM}, the first half of (iii) is
\cite[Proposition~8]{KM} (and is part (2a) of the conjecture), and the second half of (iii), stated
as an experimental observation in the remark following the conjecture, follows from
\cite[Lemma~7 and Proposition~8]{KM}. What is new is (i) for $e\ge3$ and (ii) for all $e\ge2$; we
state the other cases because the proof gives them at no cost.

The structure behind Theorems~\ref{thm:main} and~\ref{thm:sigma} is a factorisation
(Proposition~\ref{prop:structure}): each of
the series $f(z)$ in question can be written as
\[
f(z)=\bigl(T(z)+r(z)\bigr)\,\exp\!\Bigl(\nu\,\frac{z^{p+1}}{p}\Bigr)\bigl(1+p\,h(z)\bigr),
\]
where $T$ is a polynomial of degree at most $p$ with integral coefficients, $\nu$ is a unit, and
$r/p$ and $h$ lie in the ideal $J$ of Section~\ref{sec:ring}.

\subsection*{Organisation} Section~\ref{sec:ring} introduces the ring $D$ of power series with the
relevant denominators and proves the factorisation criterion. Section~\ref{sec:bky} recalls the
results of \cite{BKY} that we use and verifies the criterion for $\sigma(z)/z$ and
$\sigma(z)\sqrt{\wp(z)}$. Section~\ref{sec:romik} identifies $d(n)$ with Taylor coefficients of
$\sigma(z)\sqrt{\wp(z)}$. Section~\ref{sec:proofs} proves Theorems~\ref{thm:main} and~\ref{thm:sigma} and their
corollaries, Section~\ref{sec:uv} proves Theorem~\ref{thm:uv} and Corollary~\ref{cor:KM43}, and
Section~\ref{sec:Lvalues} expresses $d(n)^2$ as a central value of a Hecke $L$-function
(Theorem~\ref{thm:Lvalue}), and Section~\ref{sec:numerics} records numerical data and open
points.

\section{A ring of power series with divided powers}\label{sec:ring}

Let $K$ be a finite unramified extension of $\Q_p$ with ring of integers $\cO$ and valuation
$v=v_p$, so that $v(K^\times)=\Z$.

\begin{definition}
For $n\ge0$ put $m(n)=\fl{n/(p+1)}$ and
\[
\tau_n \;=\; m(n)+v_p\bigl(m(n)!\bigr)\;=\;v_p\bigl((p\,m(n))!\bigr).
\]
Let
\[
D=\Bigl\{\sum_{n\ge0}a_nz^n\in K[[z]] \;:\; v(a_n)\ge-\tau_n \text{ for all } n\Bigr\},\qquad
J=D\cap z^{p+1}K[[z]] .
\]
We write $y=z^{p+1}/p$ and $\dpow{y}{k}=y^k/k!$. Finally let
\[
\Lambda=\Bigl\{\sum_{n\ge0}a_nz^n\in K[[z]] \;:\; a_n\in\cO \text{ for } n\le p,\quad
v(a_n)\ge1-\tau_n \text{ for } n\ge p+1\Bigr\},
\]
the set of sums of a polynomial of degree at most $p$ over $\cO$ and an element of $pJ$.
\end{definition}

The numbers $\tau_n$ are those of \cite[(13)]{BKY}. Since $v(p^kk!)=\tau_{(p+1)k}$, an element of
$D$ is the same thing as a series $\sum_{k\ge0}P_k(z)\,\dpow{y}{k}$ with $P_k\in\cO[z]$ of degree at
most $p$, and this expression is unique.

\begin{lemma}\label{lem:tau}
\begin{enumerate}
\item[(a)] $\tau_{a+b}\ge\tau_a+\tau_b$ for all $a,b\ge0$, and $\tau_{n+p+1}\ge\tau_n+1$.
\item[(b)] $v_p\bigl(\fl{n/4}!\bigr)\le\tau_n-1$ for all $n\ge p+1$.
\end{enumerate}
\end{lemma}

\begin{proof}
(a) We have $m(a+b)\ge m(a)+m(b)$, and $v_p((pM)!)\ge v_p((pM_1)!)+v_p((pM_2)!)$ whenever
$M\ge M_1+M_2$, because $(pM_1+pM_2)!/((pM_1)!(pM_2)!)$ is an integer. The second assertion follows
from $\tau_{p+1}=1$.

(b) Let $k=\fl{n/4}$ and $m=m(n)\ge1$. If $k<p$ then $v_p(k!)=0\le m-1\le\tau_n-1$. Let $k\ge p$.
Since $v_p(k!)\le(k-1)/(p-1)$, it suffices to show $k+p-2\le(p-1)m$. We have
$(p+1)m\ge n-p\ge4k-p$, and
\begin{multline*}
(p-1)(4k-p)-(p+1)(k+p-2)=k(3p-5)-(2p^2-2p-2)\\
\ge p(3p-5)-(2p^2-2p-2)=(p-1)(p-2)\ge0 .
\end{multline*}
Hence $(p+1)(p-1)m\ge(p-1)(4k-p)\ge(p+1)(k+p-2)$.
\end{proof}

\begin{lemma}\label{lem:ring}
\begin{enumerate}
\item[(a)] $D$ is a subring of $K[[z]]$ containing $\cO[[z]]$, and $J$ and $pJ$ are ideals of $D$.
\item[(b)] $z^{p+1}D\subset pD$.
\item[(c)] If $x\in J$, then $\exp(px)\in1+pJ$.
\item[(d)] If $Q\in z^4\cO[[z]]$, then $\exp(Q)\in\Lambda$.
\item[(e)] $\Lambda$ is a subring of $D$.
\end{enumerate}
\end{lemma}

\begin{proof}
(a) and (b) follow from Lemma~\ref{lem:tau}(a). For (c), we have
$\exp(px)-1=\sum_{k\ge1}(p^k/k!)\,x^k$. Here $x^k\in J\cap z^{k(p+1)}K[[z]]$ because $J$ is an
ideal, so the series converges $z$-adically, and $v(p^k/k!)\ge k-\frac{k-1}{p-1}\ge1$ for $k\ge1$.

(d) The coefficient of $z^n$ in $\exp(Q)$ is $\sum_{k\ge0}[z^n]\,Q^k/k!$. Since
$Q^k\in z^{4k}\cO[[z]]$, only the terms with $k\le n/4$ contribute, so this coefficient has
valuation at least $-v_p(\fl{n/4}!)$. For $n\le p$ this is $0$, because $\fl{n/4}<p$; for
$n\ge p+1$ it is at least $1-\tau_n$ by Lemma~\ref{lem:tau}(b).

(e) Let $a,b\in\Lambda$ and write $a=A+a_1$, $b=B+b_1$ with $A,B\in\cO[z]$ of degree at most $p$
and $a_1,b_1\in pJ$. Then $ab-AB\in pJ$ because $pJ$ is an ideal of $D$, and $AB\in\cO[z]$ lies in
$\Lambda$ because $1-\tau_n\le0$ for $n\ge p+1$.
\end{proof}

\begin{proposition}[Factorisation]\label{prop:structure}
Let $L=\sum_{n\ge4}\ell_nz^n\in z^4K[[z]]$ and assume
\begin{enumerate}
\item[(a)] $\ell_n\in\cO$ for $n\le p$;
\item[(b)] $\nu:=p\,\ell_{p+1}\in\cO$;
\item[(c)] $v(\ell_n)\ge1-\tau_n$ for all $n\ge p+2$.
\end{enumerate}
Then $f=\exp(L)=\sum_{n\ge0}f_nz^n$ can be written as
\begin{equation}\label{eq:factor}
f=V\cdot\exp(\nu y)\cdot(1+ph),
\end{equation}
where $V\in\Lambda$, $V(0)=1$ and $h\in J$. In particular
$f\in D$, that is, $v(f_n)\ge-\tau_n$ for all $n$, and for all $m\ge0$ and $0\le s\le p$
\begin{equation}\label{eq:leading}
p^m\,m!\,f_{(p+1)m+s}\;\equiv\;\nu^m f_s\pmod{p\,\cO},\qquad \nu\equiv p\,f_{p+1}\pmod{p\,\cO}.
\end{equation}
More generally, if $w\in\Lambda$ and $w(0)=1$, then $fw\in D$ and \eqref{eq:leading} holds with
$f$ replaced by $fw$.
\end{proposition}

\begin{proof}
Write $L=Q+\nu y+L_1$ with $Q=\sum_{n\le p}\ell_nz^n\in z^4\cO[z]$ and
$L_1=\sum_{n\ge p+2}\ell_nz^n$, which lies in $pJ$ by (c). Then
$f=\exp(Q)\exp(\nu y)\exp(L_1)$. Lemma~\ref{lem:ring}(d) gives $V:=\exp(Q)\in\Lambda$ and
Lemma~\ref{lem:ring}(c) gives $\exp(L_1)=1+ph$ with $h\in J$. This is \eqref{eq:factor}.

Let $w\in\Lambda$ with $w(0)=1$; the case $w=1$ is the first assertion. By
Lemma~\ref{lem:ring}(e), $Vw\in\Lambda$, so $Vw=T+r$ with $T\in1+z\cO[z]$ of degree at most $p$
and $r\in pJ$. Since $\exp(\nu y)=\sum_k\nu^k\dpow{y}{k}\in D$ and $D$ is a ring, $fw\in D$, and
\[
fw-T\exp(\nu y)=r\exp(\nu y)(1+ph)+T\exp(\nu y)\,ph\in pD .
\]
Therefore $fw\equiv T(z)\exp(\nu y)=\sum_{k\ge0}\nu^kT(z)\,\dpow{y}{k}\pmod{pD}$. Comparing
coefficients of $z^{(p+1)m+s}$, and using that $T$ has degree at most $p$ and that an element of
$pD$ has coefficients of valuation at least $1-\tau_n$, we get
$p^mm!\,(fw)_{(p+1)m+s}\equiv\nu^mT_s\pmod{p\cO}$, where $T_s$ is the coefficient of $z^s$ in $T$.
The case $m=0$ gives $(fw)_s\equiv T_s$, and the case $m=1$, $s=0$ gives $p\,(fw)_{p+1}\equiv\nu$.
This proves \eqref{eq:leading} for $fw$.
\end{proof}

We isolate the two elementary inequalities used in Section~\ref{sec:bky}.

\begin{lemma}\label{lem:ineq}
Let $n\ge p+1$, $m=m(n)$ and $q=\fl{n/(p^2-1)}$.
\begin{enumerate}
\item[(a)] $m\ge2q$ and $m-q\ge\lceil m/2\rceil\ge1$.
\item[(b)] If $n\ge4p^k$ for some $k\ge1$, then $\tau_n-q\ge k+1$.
\end{enumerate}
\end{lemma}

\begin{proof}
(a) $n/(p+1)\ge(p-1)q$ and $(p-1)q$ is an integer, so $m\ge(p-1)q\ge2q$; and $m\ge1$.

(b) Since $4p\ge3(p+1)$ we have $m\ge\fl{4p^k/(p+1)}\ge3p^{k-1}$. Hence $m-q\ge\lceil m/2\rceil\ge2$,
and $v_p(m!)\ge\fl{m/p}\ge3p^{k-2}\ge k-1$ if $k\ge2$, while $v_p(m!)\ge0=k-1$ if $k=1$. So
$\tau_n-q=(m-q)+v_p(m!)\ge k+1$.
\end{proof}

\section{Weierstrass functions at supersingular primes}\label{sec:bky}

Let $E$ be the elliptic curve $y^2=4x^3-g_2x-g_3$ with $g_2,g_3\in\cO$ and good reduction, and let
$\sigma(z)$, $\zeta(z)=\sigma'(z)/\sigma(z)$, $\wp(z)=-\zeta'(z)$ be its Weierstrass functions,
regarded as formal series in $z$ as in \cite[\S1]{BKY}. This is the setting of \cite{BKY}: $p$
is an odd prime, $K$ is unramified over $\Q_p$, and $E$ is given by a Weierstrass model
$y^2=4x^3-g_2x-g_3$ over $\cO$ with good reduction. We use the following results.

\begin{theorem}[Bannai--Kobayashi--Yasuda]\label{thm:bky}
Assume that $E$ has supersingular reduction.
\begin{enumerate}
\item[(a)] \cite[Theorem~5.6]{BKY} If $\sigma(z)=\sum a_nz^n$, then $v(a_n)\ge-\tau_{n-1}$; that
is, $\sigma(z)/z\in D$.
\item[(b)] \cite[Corollary~5.3]{BKY} For $n\ge-2$ the coefficient of $z^n$ in $\wp(z)$ has
valuation at least $-(n+2)/(p^2-1)$.
\item[(c)] \cite[Corollary~5.5]{BKY} If $\zeta(z)=\sum_{n\ge-1}b_nz^n$, then
\[
v(b_n)\ge-1-\frac{n-p}{p^2-1}\quad\text{for all }n,\qquad
v(b_n)\ge-\frac{n+1}{p^2-1}\quad\text{if }p\nmid n,
\]
and equality holds in the first bound for $n=p$, that is, $v(b_p)=-1$ (see also
\cite[Proposition~4.3]{BKY}).
\end{enumerate}
\end{theorem}

Let $R=\{\sum r_nz^n\in K[[z]]:v(r_n)\ge-n/(p^2-1)\}$. This is a ring, and
$R\subset\Lambda\subset D$: the valuation $v$ is integer valued, $\fl{n/(p^2-1)}=0$ for $n\le p$,
and $\fl{n/(p^2-1)}\le m(n)-1\le\tau_n-1$ for $n\ge p+1$ by Lemma~\ref{lem:ineq}(a).

\begin{proposition}\label{prop:general}
Assume that $E$ has supersingular reduction, and let $e\in\cO$ be a root of $4x^3-g_2x-g_3$. Let
$\sigma_e(z)=\sigma(z)\sqrt{\wp(z)-e}\in1+zK[[z]]$. Then $\sigma_e\in D$.
\end{proposition}

\begin{proof}
By Theorem~\ref{thm:bky}(b), $g(z):=z^2(\wp(z)-e)$ satisfies $g-1\in R\cap zK[[z]]$. Since $p$
is odd, the binomial series $g^{1/2}=\sum_k\binom{1/2}{k}(g-1)^k$ has coefficients in $\Z_p$ and
converges $z$-adically, so $g^{1/2}\in R\subset\Lambda$. Now $\sigma_e=(\sigma(z)/z)\cdot g^{1/2}$ and $\sigma(z)/z\in D$ by
Theorem~\ref{thm:bky}(a).
\end{proof}

From now on we assume $p\equiv3\pmod4$, $g_3=0$ and $g_2\in\cO^\times$. Then $E$ has good
supersingular reduction: the discriminant is $g_2^3$ up to a power of $2$, and the Hasse invariant
is the coefficient of $x^{p-1}$ in $x^{(p-1)/2}(4x^2-g_2)^{(p-1)/2}$
\cite[Chapter~V, Theorem~4.1]{Silverman}, which vanishes because $(p-1)/2$ is odd. Moreover $\sigma(z)\in zK[[z^4]]$ and $z^2\wp(z)\in K[[z^4]]$ by
\cite[(3)]{BKY}, and $e=0$ is a root of $4x^3-g_2x$.

\begin{remark}\label{rem:three}
The hypotheses of \cite{BKY} allow $p=3$. In that case the only input beyond the bounds of
Theorem~\ref{thm:bky}(b),(c) is $v(b_3)=-1$, which can be seen directly: $b_3=-g_2/60$.
\end{remark}

\begin{proposition}\label{prop:verify}
Let $p\equiv3\pmod4$ and $g_3=0$, $g_2\in\cO^\times$.
\begin{enumerate}
\item[(1)] The series $L_\sigma=\log\bigl(\sigma(z)/z\bigr)$ lies in $z^4K[[z^4]]$ and satisfies
\textup{(a)--(c)} of Proposition~\ref{prop:structure}, with $p\,\ell_{p+1}\in\cO^\times$.
\item[(2)] The series $w=\bigl(z^2\wp(z)\bigr)^{1/2}\in1+z^4K[[z^4]]$ lies in $\Lambda$.
\end{enumerate}
Consequently $\sigma(z)/z$ and $\sigma(z)\sqrt{\wp(z)}=(\sigma(z)/z)\cdot w$ both lie in $D$ and
satisfy \eqref{eq:leading}, with the same unit $\nu=p\,\ell_{p+1}$.
\end{proposition}

\begin{proof}
(2) was shown in the proof of Proposition~\ref{prop:general} (with $e=0$), and the last assertion
follows from (1), (2) and Proposition~\ref{prop:structure}. It remains to prove (1).

Since $\frac{d}{dz}L_\sigma=\zeta(z)-1/z$, the coefficient of $z^n$
in $L_\sigma$ is $\ell_n=b_{n-1}/n$ for $n\ge2$. It vanishes unless $4\mid n$. Let $4\mid n$.

If $n<p+1$, then $p\nmid n(n-1)$ and Theorem~\ref{thm:bky}(c) gives $v(b_{n-1})\ge-n/(p^2-1)>-1$,
so $\ell_n\in\cO$. If $n=p+1$, then $v(\ell_{p+1})=v(b_p)=-1$.

Let $n\ge p+2$, and let $m,q$ be as in Lemma~\ref{lem:ineq}. Suppose first $p\nmid n-1$. Then
$v(b_{n-1})\ge-q$ by Theorem~\ref{thm:bky}(c), so $v(\ell_n)\ge-q-k$ with $k=v_p(n)$. If $k=0$ this
is at least $1-m\ge1-\tau_n$ by Lemma~\ref{lem:ineq}(a). If $k\ge1$, then $n=p^ka$ with $4\mid a$,
so $n\ge4p^k$ and Lemma~\ref{lem:ineq}(b) gives $-q-k\ge1-\tau_n$. Suppose now $n-1=pc$. Then
$p\nmid n$, and $n\equiv0$, $p\equiv3\pmod4$ force $c\equiv1\pmod4$; since $n\ne p+1$ we get
$c\ge5$, $n\ge5p+1$ and $m\ge4$. Theorem~\ref{thm:bky}(c) gives
$v(\ell_n)=v(b_{n-1})\ge-1-\fl{(n-1-p)/(p^2-1)}\ge-1-q$, and $m-q\ge\lceil m/2\rceil\ge2$ by
Lemma~\ref{lem:ineq}(a), so $v(\ell_n)\ge1-m\ge1-\tau_n$.
\end{proof}

\section{Romik's sequence and the lemniscatic curve}\label{sec:romik}

Let $\mathcal L=\varpi\,\Z[i]\subset\C$ with $\varpi=\Gamma(1/4)^2/(2\sqrt{2\pi})$; this is the period
lattice of $dx/y$ on $y^2=4x^3-4x$, with invariants $g_2=4$, $g_3=0$
\cite[(23.5.4)]{DLMF}. Let $\sigma$ and $\wp$ be its Weierstrass functions.

\begin{proposition}\label{prop:identify}
Let $\sigma(z)\sqrt{\wp(z)}=\sum_{j\ge0}f_jz^j$ be the Taylor expansion at $z=0$, normalised by
$f_0=1$. Then $f_j=0$ unless $4\mid j$, and for all $n\ge0$
\[
d(n)=2^{-n}\,(4n)!\,f_{4n}.
\]
\end{proposition}

\begin{proof}
Let $\vartheta(u)=\sum_{n\in\Z}e^{-\pi n^2+2\pi inu}$ and $\Theta(u)=e^{\pi u^2/2}\vartheta(u)$, an
entire function with $\Theta(0)=\theta_3(1)$.

\emph{Step 1.} For $a,b\in\C$ we have the identity of power series in $t$
\begin{equation}\label{eq:gauss}
\sum_{k\ge0}\frac{t^k}{k!}\,\frac{d^{2k}}{du^{2k}}e^{au^2+bu}\Big|_{u=0}
=(1-4at)^{-1/2}\exp\!\Bigl(\frac{b^2t}{1-4at}\Bigr).
\end{equation}
Indeed, the coefficient of $a^jb^{2m}t^{m+j}$ is $\frac{(2m+2j)!}{(m+j)!\,j!\,(2m)!}$ on the left
and $\frac{4^j}{m!}\binom{m+j-1/2}{j}$ on the right, and both equal
$\frac{2^j}{m!\,j!}\prod_{i=1}^j(2m+2i-1)$.

Apply \eqref{eq:gauss} with $a=\pi/2$, $b=2\pi in$ and $t=-z/(2\pi)$, and multiply by
$e^{-\pi n^2}$. The right-hand side becomes
$(1+z)^{-1/2}\exp\bigl(-\pi n^2\frac{1-z}{1+z}\bigr)$. Summing over $n\in\Z$, which is legitimate
for the Taylor coefficients at $z=0$ because the series converges uniformly on compact subsets of
$|z|<1$, and differentiating the series for $\Theta$ termwise, we obtain
\[
\sigma_3(z)=\sum_{k\ge0}\frac{(-1)^k\,\Theta^{(2k)}(0)}{(2\pi)^k\,k!}\,z^k .
\]
Comparing with \eqref{eq:romik} gives $\Theta^{(2k)}(0)=0$ for odd $k$ and
\begin{equation}\label{eq:d-theta}
d(n)=\frac{\Theta^{(4n)}(0)}{\Theta(0)\,(4\pi^2\Phi)^n}.
\end{equation}

\emph{Step 2.} Let $2\omega_1=\varpi$ and $2\omega_3=i\varpi$, so $\tau=\omega_3/\omega_1=i$ and
$q=e^{-\pi}$. By \cite[(23.6.9), (23.6.6)]{DLMF},
\begin{align*}
\sigma(z)&=2\omega_1\exp\!\Bigl(\frac{\eta_1z^2}{2\omega_1}\Bigr)
\frac{\theta_1(\pi z/(2\omega_1),q)}{\pi\,\theta_1'(0,q)},\\
\wp(z)-e_2&=\Bigl(\frac{\pi\,\theta_2(0,q)\,\theta_4(0,q)\,\theta_3(\pi z/(2\omega_1),q)}
{2\omega_1\,\theta_1(\pi z/(2\omega_1),q)}\Bigr)^{2},
\end{align*}
where $\theta_3(w,q)=1+2\sum_{n\ge1}q^{n^2}\cos(2nw)$ \cite[(20.2.3)]{DLMF},
$\eta_1=\zeta(\omega_1)$ and $e_2=\wp(\omega_2)$ with $\omega_1+\omega_2+\omega_3=0$
\cite[(23.2.1), (23.2.12), (23.3.9)]{DLMF}. Multiplying, $\sigma(z)^2(\wp(z)-e_2)$ equals
$\exp(\eta_1z^2/\omega_1)\,\theta_3(\pi z/(2\omega_1),q)^2$ times a constant, and the value $1$
at $z=0$ fixes the constant. Hence
\[
\sigma(z)\sqrt{\wp(z)-e_2}=\exp\!\Bigl(\frac{\eta_1z^2}{2\omega_1}\Bigr)
\frac{\theta_3(\pi z/(2\omega_1),q)}{\theta_3(0,q)} .
\]
In the lemniscatic case $\omega_3=i\omega_1$ we have $e_2=0$ and $\eta_1=\pi/(4\omega_1)$
\cite[(23.5.2), (23.5.3)]{DLMF}. (Indeed $i\mathcal L=\mathcal L$ gives $\wp(iz)=-\wp(z)$ and
$\zeta(iz)=-i\zeta(z)$; since $i\omega_2\equiv\omega_2\pmod{\mathcal L}$ this gives $e_2=-e_2$, and
$\eta_3=-i\eta_1$, and the Legendre relation
$\eta_1\omega_3-\eta_3\omega_1=\pi i/2$ \cite[(23.2.14)]{DLMF} gives the value of $\eta_1$.) With
$u=z/(2\omega_1)$ this shows
\[
\sigma(z)\sqrt{\wp(z)}=\frac{\Theta(u)}{\Theta(0)},\qquad u=\frac{z}{\varpi}.
\]
In particular $f_j=0$ unless $4\mid j$, and
$\Theta^{(4n)}(0)/\Theta(0)=\varpi^{4n}(4n)!\,f_{4n}$.

\emph{Step 3.} We have $\varpi^4=\Gamma(1/4)^8/(64\pi^2)$ and
$4\pi^2\Phi=\Gamma(1/4)^8/(32\pi^2)$, so $\varpi^4/(4\pi^2\Phi)=1/2$. Inserting this into
\eqref{eq:d-theta} gives the claim.
\end{proof}

\section{Proofs of the main results}\label{sec:proofs}

\begin{proof}[Proof of Theorem~\ref{thm:sigma}]
By Proposition~\ref{prop:verify}, the series $f=\sigma(z)/z$, whose coefficients are
$f_n=a_{n+1}$, lies in $D$ and satisfies \eqref{eq:leading} with $\nu\in\cO^\times$. For $s=0$ we
have $a_1=1$.
\end{proof}

\begin{proof}[Proof of Theorem~\ref{thm:main}]
Take $K=\Q_p$ and $g_2=4$. By Proposition~\ref{prop:verify}, the series
$f=\sigma(z)\sqrt{\wp(z)}=\sum f_jz^j$ lies in $D$ and satisfies \eqref{eq:leading} with
$\nu\equiv pf_{p+1}$ a unit. By
Proposition~\ref{prop:identify}, $f_{4k}=\delta(k)$. Since $4n=(p+1)m+4r$ with $0\le4r\le p-3$, the
bound $v_p(f_{4n})\ge-\tau_{4n}$ is (i), and \eqref{eq:leading} with $s=4r$ is the congruence in
(ii). Finally $\delta(r)=2^rd(r)/(4r)!$ and $(4r)!$ is prime to $p$, so $\delta(r)$ is a unit exactly when
$p\nmid d(r)$.
\end{proof}

\begin{lemma}\label{lem:floor}
For all integers $A\ge0$ we have $v_p(A!)-\tau_A\ge\fl{A/p^2}$.
\end{lemma}

\begin{proof}
Let $m=m(A)$. By Legendre's formula,
$v_p(A!)-v_p((pm)!)=\sum_{i\ge1}\bigl(\fl{A/p^i}-\fl{m/p^{i-1}}\bigr)$. Since $A\ge(p+1)m$, we have
$\fl{A/p}\ge m+\fl{m/p}$ and $\fl{A/p^i}\ge\fl{m/p^{i-1}}$ for $i\ge3$. Hence the sum is at least
$\fl{m/p}+\bigl(\fl{A/p^2}-\fl{m/p}\bigr)=\fl{A/p^2}$.
\end{proof}

\begin{proof}[Proof of Corollary~\ref{cor:KM}]
Combine Theorem~\ref{thm:main}(i) with Lemma~\ref{lem:floor} for $A=4n$. If $n\ge ep^2/4$ then
$\fl{4n/p^2}\ge e$.
\end{proof}

\begin{proof}[Proof of Corollary~\ref{cor:modp}]
If $n>(p^2-1)/4$ then $4n\ge p^2$ and Corollary~\ref{cor:KM} gives $p\mid d(n)$. For
$n=(p^2-1)/4$ we have $4n=(p+1)(p-1)$, so $m=p-1$, $r=0$ and
$v_p((4n)!)-v_p((pm)!)=(p-1)-(p-1)=0$. Since $d(0)=1$, Theorem~\ref{thm:main}(ii) gives
$v_p(d(n))=0$.
\end{proof}

\section{The sequences $u(n)$ and $v(n)$}\label{sec:uv}

Throughout this section $p\equiv3\pmod4$, except in Remark~\ref{rem:split}. Put
$\Pi_1(N)=\prod_{j=1}^N(4j-1)^2$ and $\Pi_3(N)=\prod_{j=1}^N(4j-3)^2$. Expanding the Pochhammer
symbols gives (cf.\ \cite[\S2]{KM})
\begin{equation}\label{eq:hyp}
F_1(4t^2)=\sum_{j\ge0}\frac{\Pi_3(j)}{(2j)!}\,t^{2j},\qquad
t\,F_3(4t^2)=\sum_{j\ge0}\frac{\Pi_1(j)}{(2j+1)!}\,t^{2j+1}.
\end{equation}
Let
\[
S=\Bigl\{\sum_{k\ge0}a_kt^k\in\Q_p[[t]]\;:\;v_p(a_k)\ge-\fl{\frac{k}{p+1}}\text{ for all }k\Bigr\}.
\]
Since $\fl{i/(p+1)}+\fl{j/(p+1)}\le\fl{(i+j)/(p+1)}$, this is a subring of $\Q_p[[t]]$ containing
$\Z_p[[t]]$.

\begin{lemma}\label{lem:products}
For all $j\ge0$,
\[
v_p\bigl(\Pi_3(j)\bigr)\ge v_p\bigl((2j)!\bigr)-\fl{\frac{2j}{p+1}},\qquad
v_p\bigl(\Pi_1(j)\bigr)\ge v_p\bigl((2j+1)!\bigr)-\fl{\frac{2j+1}{p+1}} .
\]
That is, both series in \eqref{eq:hyp} lie in $S$.
\end{lemma}

\begin{proof}
Fix $i\ge1$, put $P=p^i$ and write $j=aP+b$ with $0\le b<P$. For $r\in\{1,3\}$ let
$c_r\in\{1,\dots,P\}$ be the solution of $4c_r\equiv r\pmod P$, so that
$\#\{1\le l\le j:P\mid 4l-r\}=a+[b\ge c_r]$, where $[\,\cdot\,]$ is $1$ if the condition holds and
$0$ otherwise. If $i$ is odd then $P\equiv3\pmod4$, $c_3=3(P+1)/4$ and $c_1=(P+1)/4$; if $i$ is
even then $P\equiv1\pmod4$, $c_3=(P+3)/4$ and $c_1=(3P+1)/4$. Moreover
$\fl{2j/P}=2a+[b\ge(P+1)/2]$ and $\fl{(2j+1)/P}=2a+[b\ge(P-1)/2]$. Hence
\begin{align*}
2\,\#\{l\le j:P\mid4l-3\}-\fl{2j/P}&=2[b\ge c_3]-[b\ge(P+1)/2],\\
2\,\#\{l\le j:P\mid4l-1\}-\fl{(2j+1)/P}&=2[b\ge c_1]-[b\ge(P-1)/2].
\end{align*}
The first expression is at least $-1$, and it is non-negative unless $i$ is odd and
$b\ge(P+1)/2$, which forces $p^i<2j$. The second is at least $-1$, and it is non-negative unless
$i$ is even and $b\ge(P-1)/2$, which forces $p^i\le2j+1$. Summing over $i$ and using Legendre's
formula together with $v_p(\Pi_r(j))=2\sum_{i\ge1}\#\{l\le j:p^i\mid4l-r\}$,
\begin{align*}
v_p(\Pi_3(j))-v_p((2j)!)&\ge-\#\{i\text{ odd}:p^i<2j\},\\
v_p(\Pi_1(j))-v_p((2j+1)!)&\ge-\#\{i\ge2\text{ even}:p^i\le2j+1\}.
\end{align*}
If the first set has $s\ge1$ elements, then $2j\ge p^{2s-1}+1$ and
$\fl{2j/(p+1)}\ge(p^{2s-1}+1)/(p+1)=1+\sum_{h=1}^{s-1}p^{2h-1}(p-1)\ge s$. If the second set has
$s\ge1$ elements, then $2j+1\ge p^{2s}$ and
$\fl{(2j+1)/(p+1)}\ge(p^{2s}-1)/(p+1)=(p-1)\sum_{h=0}^{s-1}p^{2h}\ge s$.
\end{proof}

\begin{proof}[Proof of Theorem~\ref{thm:uv}]
By Lemma~\ref{lem:products}, $F_1(4t^2)\in S$ and $tF_3(4t^2)\in S$. The ring $S$ is defined by
conditions on the individual coefficients, so it is closed under $t$-adically convergent sums.
Since $F_1(4t^2)\in1+t^2\Q_p[[t]]$, the geometric series gives $F_1(4t^2)^{-1}\in S$, and, $p$
being odd, the binomial series gives $F_1(4t^2)^{1/2}\in S$ for the square root with constant
term $1$. Hence both series in \eqref{eq:uv} lie in $S$,
which is the assertion.
\end{proof}

\begin{lemma}\label{lem:B}
For $k\ge0$ let $B(k)=v_p(k!)-\fl{k/(p+1)}$, and write $k=(p+1)a+b$ with $0\le b\le p$. Then
\[
B(k)=\fl{\frac{a+b}{p}}+\sum_{i\ge2}\fl{\frac{k}{p^i}} .
\]
Consequently:
\begin{enumerate}
\item[(a)] $B(k)\ge2e-1$ for all $k\ge ep^2$, $e\ge1$;
\item[(b)] $B(k)\ge2e$ for all odd $k\ge ep^2+(e-2)p+1$, $e\ge2$;
\item[(c)] $B(k)\ge2$ for all $k\ge p^2$, and $B(k)\ge1$ for all odd $k\ge p^2-p+1$.
\end{enumerate}
\end{lemma}

\begin{proof}
The formula follows from Legendre's formula and $\fl{k/p}=\fl{(pa+a+b)/p}=a+\fl{(a+b)/p}$. Note
that $a+b\ge k/(p+1)$ and $k=pa+(a+b)$.

(a) Let $k\ge ep^2$. Then $a+b\ge ep^2/(p+1)\ge e(p-1)$, so
$\fl{(a+b)/p}\ge\fl{e(p-1)/p}=e-\lceil e/p\rceil$, and
$\sum_{i\ge2}\fl{k/p^i}\ge e+\fl{e/p}$. Hence $B(k)\ge2e+\fl{e/p}-\lceil e/p\rceil\ge2e-1$.

(b) Let $k\ge ep^2+(e-2)p+1=(p+1)(ep-2)+3$ be odd. Then $a+b>ep-2$, so $a+b\ge ep-1$. If
$a+b=ep-1$, then $pa=k-(ep-1)\ge p(ep-2)+2$, so $a\ge ep-1$, hence $b=0$ and $k=(p+1)(ep-1)$ is
even, a contradiction. Therefore $a+b\ge ep$ and $\fl{(a+b)/p}\ge e$. Since $k\ge ep^2$, the sum
over $i\ge2$ is at least $e$.

(c) Let $k\ge p^2$. Then $\fl{k/p^2}\ge1$, and $a+b\ge p^2/(p+1)>p-1$ gives $\fl{(a+b)/p}\ge1$. Let
$k\ge p^2-p+1=(p+1)(p-2)+3$ be odd. Then $a+b\ge p-1$, and $a+b=p-1$ would give
$pa=k-(p-1)\ge p(p-2)+2$, so $a\ge p-1$, $b=0$ and $k=p^2-1$ even. Hence $a+b\ge p$.
\end{proof}

\begin{proof}[Proof of Corollary~\ref{cor:KM43}]
By Theorem~\ref{thm:uv}, $v_p(u(n))\ge B(2n+1)$. If $n\ge\lceil(ep^2-1)/2\rceil$ then
$2n+1\ge ep^2$; if $n\ge(ep^2+(e-2)p)/2$ then $2n+1\ge ep^2+(e-2)p+1$; if $n\ge(p^2-1)/2$ then
$2n+1\ge p^2$; if $n\ge(p^2-p)/2$ then $2n+1\ge p^2-p+1$. Apply Lemma~\ref{lem:B}.
\end{proof}

\begin{remark}\label{rem:v}
For $v(n)$, Theorem~\ref{thm:uv} gives $v(n)\equiv0\pmod{p^e}$ whenever $B(2n)\ge e$. By
Lemma~\ref{lem:B}(a),(c) this holds for all $n\ge\lceil(e-1)p^2/2\rceil$, $e\ge2$, which is
\cite[Theorem~17]{KM}; in our computations the thresholds obtained from $B$ are smaller for
$e\ge4$. They give \cite[Conjecture~44]{KM} for some pairs $(p,e)$ but not for others: for $p=7$
and $e=1$ the bound gives $n\ge25$, while the conjecture predicts $n\ge18$. Equality in the bound
is less frequent for $v(n)$ than for $u(n)$. In our data ($n\le1200$, $p\in\{3,7,11,19,23\}$) it
occurs for $v(n)$ only when the binomial coefficient $\binom{1/2}{m}$, $m=\fl{2n/(p+1)}$, of the
expansion of $F_1(4t^2)^{1/2}$ is a $p$-adic unit, whereas for $u(n)$, where only a geometric series
is involved, it occurs for more than $80\%$ of all $n$.
\end{remark}

\begin{remark}\label{rem:sporadic}
The bound for $u(n)$ is attained at $n=(3p^2+p-2)/2$ for $p\in\{7,11,19,23\}$, where
$v_p(u(n))=5$. Hence the second of the two sporadic improvements suggested in the remark following
\cite[Conjecture~43]{KM}, namely $u(n)\equiv0\pmod{p^6}$ for $n\ge(3p^2-1)/2$, does not hold for
these primes. Similarly $v_3(v(10))=3$, so the threshold $\lceil ep^2/4\rceil$ of
\cite[Conjecture~44(2)]{KM} fails for $(p,e)=(3,4)$.
\end{remark}

\begin{remark}\label{rem:split}
For $p\equiv1\pmod4$ the proof of Lemma~\ref{lem:products} shows that $F_1(4t^2)\in\Z_p[[t]]$
(now every $P=p^i$ is $\equiv1\pmod4$), hence $v_p(v(n))\ge v_p((2n)!)$ for all $n$; compare
\cite[Theorem~19]{KM}.
\end{remark}

\section{Central values of Hecke $L$-functions}\label{sec:Lvalues}

Let $\psi$ be the Hecke character of $\Q(i)$ attached to the elliptic curve $y^2=x^3-x$ of
conductor $32$. Every ideal of $\Z[i]$ prime to $2$ has a unique generator
$\alpha\equiv1\pmod{2+2i}$, and $\psi((\alpha))=\alpha$ for such $\alpha$
\cite[Chapter~18, \S\S4 and~6]{IR}. Explicitly, $\alpha=a+bi\equiv1\pmod{2+2i}$ means that $a$
is odd, $b$ is even and $a+b\equiv1\pmod4$. For odd $j\ge1$ the function
\[
L(\psi^j,s)=\sum_{\alpha\equiv1\ (2+2i)}\frac{\alpha^j}{N(\alpha)^s}\qquad(\operatorname{Re}s>\tfrac j2+1)
\]
continues to an entire function of $s$ \cite{Hecke18}, with centre $s=(j+1)/2$. Let $\varpi$ be as in
Section~\ref{sec:romik}.

\begin{theorem}\label{thm:Lvalue}
For all $n\ge0$,
\[
L\bigl(\psi^{8n+1},\,4n+1\bigr)\;=\;\frac{\varpi^{8n+1}}{\pi^{4n}\,(4n)!}\;4^{\,n-1}\,d(n)^2
\qquad\text{and}\qquad L\bigl(\psi^{8n+5},\,4n+3\bigr)=0 .
\]
\end{theorem}

For $n=0$ this reads $L(\psi,1)=\varpi/4$.

\begin{corollary}\label{cor:Lvalue}
Put $L^{\mathrm{alg}}(n)=\pi^{4n}(4n)!\,L(\psi^{8n+1},4n+1)/\varpi^{8n+1}\in\Q$. Then
$v_p(L^{\mathrm{alg}}(n))=2\,v_p(d(n))$ for every odd prime $p$. If $p\equiv3\pmod4$, then with
$m$ and $r$ as in Theorem~\ref{thm:main},
\[
v_p\bigl(L^{\mathrm{alg}}(n)\bigr)\;\ge\;2\Bigl(v_p\bigl((4n)!\bigr)-v_p\bigl((pm)!\bigr)\Bigr),
\]
with equality if and only if $p\nmid d(r)$.
\end{corollary}

The primes $p\equiv3\pmod4$ are those inert in $\Q(i)$. In Corollary~\ref{cor:Lvalue} the
conductor of the character is prime to $p$ and its infinity type varies. For a prime $p$ inert
in an imaginary quadratic field $K$, Burungale, Kobayashi and Ota \cite[Theorem~5.3]{BKO}
determined $v_p\bigl(L(\varphi\chi,1)/\Omega\bigr)$, where $\Omega$ is a CM period, $\varphi$ is a
self-dual Hecke character of $K$ of infinity type $(1,0)$ with
$p\nmid6h_K\operatorname{cond}\varphi$, and $\chi$ is an anticyclotomic character of order $p^j$
with $j$ sufficiently large and $(-1)^{j-1}$ equal to the root number of $\varphi$; see also
\cite{BHKO}, where the case of general infinity type is mentioned as a natural problem.

Theorem~\ref{thm:Lvalue} is of the type established by Rodriguez Villegas and Zagier \cite{VZ}
for a Hecke character of $\Q(\sqrt{-7})$, and its proof follows the proof of
\cite[Theorem~2.10]{Nomoto}, where Nomoto treats the curve $y^2=x^3+x$ and the theta function
$\theta_2$. The connection between Taylor coefficients at CM points and $L$-values is mentioned
as background in \cite{Romik,GMR,KM}, but we have not found Theorem~\ref{thm:Lvalue} stated there.
For Taylor expansions of modular forms at CM points in general see
\cite[Sections~5.1 and~6.3]{Zagier}.

\subsection*{Notation} For $z=x+iy$ in the upper half plane let $D=\frac1{2\pi i}\frac{d}{dz}$,
$\partial_k=D-\frac{k}{4\pi y}$ and
$\partial_k^{(h)}=\partial_{k+2h-2}\circ\dots\circ\partial_{k+2}\circ\partial_k$. For
$\gamma\in\mathrm{GL}_2^+(\mathbb R)$ one has
$\partial_k(f|_k\gamma)=(\partial_kf)|_{k+2}\gamma$. Let
\[
\vartheta_K(z)=\sum_{\lambda\in\Z[i]}e^{2\pi iN(\lambda)z}=\theta_3(-2iz)^2 ,
\]
let $L_p$ be the Laguerre polynomial $L_p(t)=\sum_{j=0}^p\binom pj\frac{(-t)^j}{j!}$, and let
$H_p$ be the Hermite polynomial, $e^{2xt-t^2}=\sum_pH_p(x)t^p/p!$. Throughout the proof $p=2h$ is
an even integer $\ge0$, and
\[
\Lambda_s(\lambda)=\frac{\bar\lambda^{\,2p+1}}{|\lambda|^{2p+2+2s}}\qquad(\lambda\in\Z[i],\ \lambda\ne0).
\]

\begin{proof}[Proof of Theorem~\ref{thm:Lvalue}]
\emph{Step 1: the $L$-value as a value of Eisenstein series.} Let $\epsilon$ be the non-trivial
Dirichlet character modulo $4$ and, for $\operatorname{Re}s>\frac12$,
\[
E_p(z,s)=\frac12\sum_{(n,m)\ne(0,0)}\epsilon(n)\,\frac{(n+4m\bar z)^{2p+1}}{|n+4mz|^{2p+2+2s}} .
\]
By Hecke's method \cite{Hecke27}, $E_p(z,s)$ continues analytically to $s=0$ (see Step~2); put $E_p(z)=E_p(z,0)$. We have
$E_p(i,s)=\frac12\sum\epsilon(a)\Lambda_s(a+bi)$ over odd $a$ and $b\equiv0\pmod4$, and
$E_p(i/2,s)=\frac12\sum\epsilon(a)\Lambda_s(a+bi)$ over odd $a$ and even $b$. Let $P_0(s)$ and
$P_2(s)$ be the sums of $\Lambda_s(\alpha)$ over the $\alpha=a+bi\equiv1\pmod{2+2i}$ with
$b\equiv0$, respectively $b\equiv2\pmod4$. If $a$ is odd and $b\equiv0\pmod4$, then exactly one of
$\pm(a+bi)$ is $\equiv1\pmod{2+2i}$, namely the one with real part $\equiv1\pmod4$; if
$b\equiv2\pmod4$ it is the one with real part $\equiv3\pmod4$. Since
$\Lambda_s(-\lambda)=-\Lambda_s(\lambda)$, this gives
\[
E_p(i,s)=P_0(s),\qquad E_p(i/2,s)=P_0(s)-P_2(s).
\]
The set of $\alpha\equiv1\pmod{2+2i}$ is stable under complex conjugation, so
$L(\psi^{2p+1},p+1+s)=P_0(s)+P_2(s)$, and at $s=0$
\begin{equation}\label{eq:L-eis}
L(\psi^{2p+1},p+1)=2E_p(i)-E_p(i/2).
\end{equation}

\emph{Step 2: Eisenstein series and $\vartheta_K$.} For $p=0$, Hecke's identity
$E_0(z)=\frac\pi4\vartheta_K(z)$ holds (see \cite[proof of Proposition~2.7]{Nomoto}; it is
equivalent to Jacobi's formula $r_2(N)=4\sum_{d\mid N}\epsilon(d)$ for the number of
representations of $N$ as a sum of two squares \cite[Theorem~278]{HW}). Applying $\partial_1^{(p)}$, by
\cite[(16)]{VZ},
\begin{equation}\label{eq:eis-theta}
\frac\pi4\,\partial_1^{(p)}\vartheta_K(z)=\frac{p!}{(4\pi y)^p}\,E_p(z).
\end{equation}
Here the differentiation is carried out for $\operatorname{Re}s>\frac12$ and then continued to
$s=0$: writing $\rho=(n+4m\bar z)/(n+4mz)$ and $\phi_s=(n+4mz)^{-1}|n+4mz|^{-2s}$, one finds
$\partial_{1+2j}\bigl(y^{-j}\rho^i\phi_s\bigr)=(4\pi)^{-1}y^{-j-1}\bigl((s+i-j)\rho^i\phi_s
-(1+i+s)\rho^{i+1}\phi_s\bigr)$, so that
$\partial_1^{(p)}\phi_s=(4\pi y)^{-p}\sum_{i\le p}c_i(s)\rho^i\phi_s$ with polynomials $c_i$
satisfying $c_i(0)=0$ for $i<p$ and $c_p(0)=p!$ (recall that $p$ is even).

It remains to justify the continuation to $s=0$ and its compatibility with differentiation. Fix
$m$ and $i$ and write $F(n)=\rho^i\phi_s$. Since $\epsilon(n+2)=-\epsilon(n)$, for
$\operatorname{Re}s>\frac12$
\[
\sum_{n}\epsilon(n)F(n)=\frac14\sum_{n}\epsilon(n)\bigl(F(n)-2F(n+2)+F(n+4)\bigr).
\]
The second difference on the right is $O\bigl(|n+4mz|^{-3-2\operatorname{Re}s}\bigr)$, locally
uniformly in $z$, and the same estimate holds after any fixed number of derivatives with respect
to $x$ and $y$. Hence the double series over $(n,m)$ obtained in this way, and its derivatives,
converge absolutely and locally uniformly for $\operatorname{Re}s>-\frac12$. This continues
$\sum\epsilon(n)\rho^i\phi_s$ to a neighbourhood of $s=0$ and shows that the continuation commutes
with differentiation in $z$; compare the remark following \cite[(22)]{VZ}.

\emph{Step 3: the Fricke involution.} From $\theta_3(1/x)=\sqrt x\,\theta_3(x)$ we get
$\vartheta_K(-1/(4z))=-2iz\,\vartheta_K(z)$, that is, $\vartheta_K|_1\gamma=-i\,\vartheta_K$ for
$\gamma=\bigl(\begin{smallmatrix}0&-1\\4&0\end{smallmatrix}\bigr)$. Hence
\[
(2z)^{-2p-1}\,\bigl(\partial_1^{(p)}\vartheta_K\bigr)(-1/(4z))=-i\,\partial_1^{(p)}\vartheta_K(z),
\]
and at $z=i$, since $p$ is even,
\begin{equation}\label{eq:fricke}
\partial_1^{(p)}\vartheta_K(i/4)=2^{2p+1}\,\partial_1^{(p)}\vartheta_K(i).
\end{equation}
By \eqref{eq:L-eis}, \eqref{eq:eis-theta} and \eqref{eq:fricke},
\begin{equation}\label{eq:L-A}
L(\psi^{2p+1},p+1)=\frac{\pi^{p+1}}{8\,p!}\,A_p,
\end{equation}
where
$A_p:=\partial_1^{(p)}\vartheta_K(i/4)-2^{p+1}\partial_1^{(p)}\vartheta_K(i/2)
+2^{2p+1}\partial_1^{(p)}\vartheta_K(i)$.

\emph{Step 4: Laguerre sums and the factorisation formula.} By \cite[(9)]{VZ}, for even $p$,
\[
\partial_1^{(p)}\vartheta_K(z)=\frac{p!}{(4\pi y)^p}\sum_{\lambda\in\Z[i]}L_p\bigl(4\pi N(\lambda)y\bigr)
e^{2\pi iN(\lambda)z}.
\]
For $\lambda=n+mi$ the number $nm$ is odd exactly when $\lambda\in(1+i)\Z[i]\setminus2\Z[i]$.
Evaluating the last formula at $z=i/4$, $i/2$ and $i$ therefore gives
\[
A_p=\frac{p!}{\pi^p}\sum_{n,m\in\Z}(-1)^{nm}\,L_p\bigl(\pi(n^2+m^2)\bigr)\,e^{-\pi(n^2+m^2)/2}.
\]
The factorisation formula of Rodriguez Villegas and Zagier \cite[(23)]{VZ} with $a=1$,
$\alpha=0$, $\mu=\nu=0$ and $z=i$ states that the
right-hand side equals $\sqrt2\,(2\pi)^{-p}\,\eta_p^2$, where
\[
\eta_p=\sum_{n\in\Z}H_p\bigl(n\sqrt{2\pi}\bigr)\,e^{-\pi n^2}.
\]

\emph{Step 5: Hermite sums and $d(n)$.} Let $\Theta(u)=e^{\pi u^2/2}\sum_ne^{-\pi n^2+2\pi inu}$ be
as in the proof of Proposition~\ref{prop:identify}. Since
$e^{\pi u^2/2+2\pi inu}=e^{2xt-t^2}$ with $t=i\sqrt{\pi/2}\,u$ and $x=n\sqrt{2\pi}$, we get
$\Theta^{(p)}(0)=(-\pi/2)^{p/2}\,\eta_p$, hence $\eta_p^2=(2/\pi)^p\,\Theta^{(p)}(0)^2$. Together
with \eqref{eq:L-A} and Step~4,
\begin{equation}\label{eq:L-Theta}
L(\psi^{2p+1},p+1)=\frac{\sqrt2}{8\,p!\,\pi^{p-1}}\;\Theta^{(p)}(0)^2\qquad(p=2h\ge0).
\end{equation}
In the proof of Proposition~\ref{prop:identify} we showed $\Theta^{(2k)}(0)=0$ for odd $k$, which
gives $L(\psi^{8n+5},4n+3)=0$ (the case $p=4n+2$), and
$\Theta^{(4n)}(0)=\theta_3(1)\,d(n)\,(4\pi^2\Phi)^n$ with $4\pi^2\Phi=2\varpi^4$. Finally
$\theta_3(1)=\Gamma(1/4)/(\sqrt2\,\pi^{3/4})$ \cite[Section~1]{Romik}, so
$\theta_3(1)^2=\sqrt2\,\varpi/\pi$. Inserting this into \eqref{eq:L-Theta} with $p=4n$ gives
\[
L(\psi^{8n+1},4n+1)=\frac{\sqrt2}{8\,(4n)!\,\pi^{4n-1}}\cdot\frac{\sqrt2\,\varpi}{\pi}\cdot
4^n\varpi^{8n}\,d(n)^2=\frac{\varpi^{8n+1}}{\pi^{4n}\,(4n)!}\,4^{n-1}\,d(n)^2 .\qedhere
\]
\end{proof}

\begin{proof}[Proof of Corollary~\ref{cor:Lvalue}]
By Theorem~\ref{thm:Lvalue}, $L^{\mathrm{alg}}(n)=4^{n-1}d(n)^2$. The rest is
Theorem~\ref{thm:main}.
\end{proof}

\begin{remark}\label{rem:Lcheck}
The following parts of the proof were checked numerically to $60$ decimal digits for
$p=2h\le20$ (script \texttt{l\_value\_proof\_check.py} in the directory \texttt{romik/}): the
case of the factorisation formula used in Step~4; the expression of $A_p$ as a Laguerre sum;
the combination of \eqref{eq:L-eis} and \eqref{eq:eis-theta}, with $L(\psi^{2p+1},p+1)$ computed
independently; \eqref{eq:fricke}; \eqref{eq:L-A} and \eqref{eq:L-Theta}; and the formula for
$\Theta^{(4n)}(0)$. The series $E_p(z,s)$ for $s\ne0$, the sums $P_0(s)$ and $P_2(s)$ and the
recursion in Step~2 were not checked numerically. The statement of
Theorem~\ref{thm:Lvalue} was checked for $n\le8$ against an independent computation of the
central values:
the Dirichlet coefficients of $L(\psi^j,s)$ were obtained from the numbers of points of
$y^2=x^3-x$ modulo primes, the functional equation was checked numerically, and the central
value was computed from the standard rapidly convergent series (script
\texttt{hecke\_l\_check.py}). For the curves $y^2=x^3+x$ and $y^2=x^3+2x$ and the exponents
$4n+1$, $n\le16$, the normalised central values computed in the same way agree to $60$ digits with
squares of rational numbers of small height; for $y^2=x^3+x$ this is explained by
\cite[Theorem~2.10]{Nomoto}.
\end{remark}

\section{Numerical data and remarks}\label{sec:numerics}

\subsection{Data} We computed $d(n)$ for $n\le800$ in two ways, from an integer recursion for the
derivatives of $\theta_3$ at $\tau=i$ and from Proposition~\ref{prop:identify} with exact
rational arithmetic; the two computations agree. For $n\le200$ the values also satisfy Romik's
recursion $v(n)=\sum_{k\le n}r(n,k)\,d(k)$ \cite[Theorem~7]{Romik}, \cite[(2.8)]{KM}, which
involves neither theta functions nor Weierstrass functions. We checked
Theorem~\ref{thm:main} for the $87$ primes $p\equiv3\pmod4$ below $1000$. For $68$ of them
equality holds in (i) for every $n\le800$. For the remaining $19$ the classes $r$ with
$p\mid d(r)$, $r<(p+1)/4$, are
\begin{gather*}
67\colon 7;\quad 71\colon 5,8,14,17;\quad 127\colon 7;\quad 239\colon 17;\quad 251\colon 62;\quad
283\colon 4;\\
347\colon 71;\quad 379\colon 8;\quad
419\colon 61,70,96;\quad 439\colon 60;\quad 499\colon 36;\quad 503\colon 28;\\
599\colon 133;\quad 607\colon 114;\quad 659\colon 96,151;\quad
743\colon 31,181;\quad 859\colon 84;\\
919\colon 184;\quad 967\colon 94,
\end{gather*}
and the $n\le800$ at which (i) is strict are exactly those in these classes. The congruence
\eqref{eq:leading} was checked
for $\sigma(z)/z$ and $\sigma(z)\sqrt{\wp(z)}$ up to degree $1200$ for $p\in\{3,7,11,19,23,31,43\}$.

For the sequences $u(n)$ and $v(n)$, Theorem~\ref{thm:uv}, Lemma~\ref{lem:products} (degrees up to
$12000$), Lemma~\ref{lem:B} ($e\le60$) and Remark~\ref{rem:split} were checked on the values for
$n\le1200$ and all primes below $300$.

\subsection{The excess} When $p\mid d(r)$, the excess
$v_p(d(n))-v_p((4n)!)+v_p((pm)!)$ equals $v_p(d(r))$ at all data points with $p\notin\{67,71\}$.
It is $2$ instead of $1$ at $(p,n)=(67,517),(71,53),(71,266),(71,692)$. We have no description of
the excess.

\subsection{Thresholds} Theorem~\ref{thm:main}(i) bounds from above the least $T$ with
$d(n)\equiv0\pmod{p^e}$ for all $n\ge T$, and by (ii) the bound is sharp whenever the last $n$ with
$v_p((4n)!)-\tau_{4n}<e$ has $p\nmid d(r)$. The threshold $\lceil ep^2/4\rceil$ of Corollary~\ref{cor:KM} is not
always optimal: for $p=3$, $e=4$ the lower bound of Theorem~\ref{thm:main}(i) is at least $4$ for
all $n\ge7$, while $\lceil ep^2/4\rceil=9$.

\subsection{Other curves} For the curve $y^2=4x^3-g_3$ and primes $p\equiv2\pmod3$ the series
$\sigma(z)/z$ lies in $K[[z^6]]$, and the analogue of Theorem~\ref{thm:sigma} holds numerically up
to degree $1200$ for $p\in\{5,11,17,23,29,41\}$; the proof should go through with $4$ replaced by
$6$, but we have not written it out. For a general curve with
supersingular reduction the argument breaks down at degree $2p$: there $L_\sigma$ is only a series
in $z^2$, so Lemma~\ref{lem:ring}(d) does not apply, and the coefficient $b_{2p-1}/(2p)$ of
$L_\sigma$ may have valuation $-1$, whereas hypothesis (c) of Proposition~\ref{prop:structure}
requires valuation at least $1-\tau_{2p}=0$. We do not know the equality pattern in that case.

\subsection{Relation to earlier work} The slope $p/(p^2-1)$ of $\tau_n$ is the valuation of the
$p$-adic period of the formal group of a supersingular curve, and the radius of convergence of the
$p$-adic sigma function \cite{MT} was determined in \cite{BKY}; a numerical example is discussed at the end
of the arXiv version of \cite{BK}. The coefficients of the lemniscatic $\wp$-function were
studied by Hurwitz \cite{Hurwitz}, and divisibility of Bernoulli--Hurwitz numbers at supersingular
primes goes back to Katz \cite{Katz}. Congruences for Taylor coefficients of modular forms at CM
points were studied in \cite{LarsonSmith}, and the behaviour of $d(n)$ modulo primes
$p\equiv1\pmod4$ in \cite{GMR,Wakhare}. The proof in \cite{KimLee} uses quasimodular forms and
the proof in \cite{KM} uses Romik's recursion; neither uses formal groups.

\subsection{Formal verification} The following parts have been formalised in Lean~4
\cite{Lean4} with Mathlib \cite{Mathlib} (directory \texttt{lean/} of the accompanying repository
\url{https://github.com/felix5572/romik-sequence-valuations}). All
formal statements depend only on the standard axioms \texttt{propext}, \texttt{Classical.choice}
and \texttt{Quot.sound}.

\begin{enumerate}
\item[(1)] \emph{The sequences $u(n)$ and $v(n)$: complete.} The files \texttt{Conj43.lean} and
\texttt{VSeq.lean} define $u(n)$ and $v(n)$ by the recurrences of \cite[\S2]{KM},
\begin{align*}
u(n)&=\Pi_1(n)-\sum_{m<n}\binom{2n+1}{2m+1}\Pi_3(n-m)\,u(m),\\
v(n)&=2^{n-1}\Pi_3(n)-\frac12\sum_{0<m<n}\binom{2n}{2m}v(m)\,v(n-m),
\end{align*}
and prove from them Theorem~\ref{thm:uv} and parts (i) and (ii) of Corollary~\ref{cor:KM43}, with
no step left on paper. There Lemma~\ref{lem:products} is replaced by the inequalities
$v_p\bigl(\prod_{l<j}(4l+a)\bigr)\ge v_p(j!)$ for $a\in\{1,3\}$ and
$v_p(n!)\le2v_p(\lfloor n/2\rfloor!)+\lfloor\log_pn\rfloor$, and the power series argument by an
induction on the recurrences. Lemma~\ref{lem:B} is in \texttt{Threshold.lean}.

\item[(2)] \emph{Section~\ref{sec:ring}, for $K=\Q_p$ and series with rational coefficients.}
Lemmas~\ref{lem:tau} and~\ref{lem:ineq}, the coefficient estimates of Lemma~\ref{lem:ring} that
are used later (for example, Lemma~\ref{lem:ring}(c) in the form: if $x\in D$ and $x(0)=0$, then
$\exp(px)-1\in pD$), and the conclusions of Proposition~\ref{prop:structure} that are used later: $fw\in D$ and the congruence
\eqref{eq:leading} for $fw$ (theorems \texttt{factorisation} and \texttt{factorisation\_mul} in
\texttt{Exp.lean}), where $f=\exp(L)$ is characterised by $f'=L'f$ and $f(0)=1$. The factorisation
\eqref{eq:factor} is formalised in the form $f=V\exp(\nu y)(1+r)$ with $V\in\Lambda$ and $r\in pD$
(theorem \texttt{exists\_factor}); the finer statements $V(0)=1$ and $r\in pJ$ are not part of
the formal statement.

\item[(3)] \emph{Theorems~\ref{thm:sigma} and~\ref{thm:main} and Corollary~\ref{cor:KM}, assuming
the cited inputs.} Let $b_n$ be rational numbers, vanishing unless $n\equiv3\pmod4$, that satisfy
the bounds of Theorem~\ref{thm:bky}(c), and let $f$, $w$ be power series with
\[
f(0)=1,\quad f'=\Bigl(\sum_nb_nz^n\Bigr)f,\qquad w(0)=1,\quad
w^2=1-z^2\Bigl(\sum_nb_nz^n\Bigr)' .
\]
For the $b_n$ of $\zeta(z)=z^{-1}+\sum_nb_nz^n$ these are $f=\sigma(z)/z$ and
$w=(z^2\wp(z))^{1/2}$, since $\wp=-\zeta'$. The theorems \texttt{sigma\_congruence}
(\texttt{Sigma.lean}) and \texttt{theta\_congruence} (\texttt{Theta.lean}) state that $f$ and $fw$
lie in $D$ and satisfy \eqref{eq:leading} with $\nu=p\,b_p/(p+1)$; they do not assert that $\nu$
is a unit. The theorem
\texttt{romik\_main} assumes in addition $v_p(b_p)=-1$ and a sequence $d$ with
$d(k)=(4k)!\,2^{-k}\,[z^{4k}](fw)$ for all $k$, which is the statement of
Proposition~\ref{prop:identify}, and concludes Theorem~\ref{thm:main}(i), the criterion for
equality in Theorem~\ref{thm:main}(ii) and Corollary~\ref{cor:KM}. The bound for the coefficients
of $\wp$ in Theorem~\ref{thm:bky}(b) is not assumed; what is needed of it is deduced from
Theorem~\ref{thm:bky}(c).
\end{enumerate}

Not formalised: the results quoted from \cite{BKY} as Theorem~\ref{thm:bky}(c) and
Proposition~\ref{prop:identify}, which enter (3) as hypotheses, and all of
Section~\ref{sec:Lvalues}. The script
\texttt{lean\_hypotheses\_check.py} in the directory \texttt{romik/} checks every hypothesis of
\texttt{romik\_main} on the actual coefficients, to degree $1200$ in $z$ and for all primes
$p\equiv3\pmod4$ below $1200$.

\section*{Acknowledgements and use of AI}

This work was prompted by the proof of the Riemann hypothesis announced by Alp\"oge, working
with Claude \cite{Alpoge}. The formula in Theorem~\ref{thm:main} was found numerically in the
course of reading an expository companion to that proof \cite{Companion}. We have not
studied the Lean source of that proof, and the proofs given here do not depend on it.

The author chose the direction of the project and its workflow: numerical experiment, literature
check, written proof, and formalisation of the elementary part. The numerical experiments, the
literature search, the proofs, the text of this paper and the Lean formalisation were produced by
Claude (Anthropic) under the author's direction. The proofs were reviewed by a second AI system,
GPT-6.2 Astra, and the author looked through the key parts.

\begin{thebibliography}{99}

\bibitem{Alpoge} L.~Alp\"oge, announcement of a proof of the Riemann hypothesis (untitled, four
pages), 20 November 2026, \url{https://alpo.ge/remark.pdf}.

\bibitem{Companion} Dirichlet $L$-functions and the critical line, unsigned expository companion to
the Lean development of \cite{Alpoge}, 23 pages, November 2026, made available by P.~Chojecki at
\url{https://www.ulam.ai/research/rh-short.pdf}.

\bibitem{BK} K.~Bannai and S.~Kobayashi, Algebraic theta functions and the $p$-adic interpolation
of Eisenstein--Kronecker numbers, \emph{Duke Math.\ J.} \textbf{153} (2010), 229--295.

\bibitem{BKY} K.~Bannai, S.~Kobayashi and S.~Yasuda, The radius of convergence of the $p$-adic
sigma function, \emph{Math.\ Z.} \textbf{286} (2017), 751--781.

\bibitem{BHKO} A.~Burungale, W.~He, S.~Kobayashi and K.~Ota, Hecke $L$-values, definite Shimura
sets and mod $\ell$ non-vanishing, preprint, arXiv:2408.13932.

\bibitem{BKO} A.~Burungale, S.~Kobayashi and K.~Ota, The $p$-adic valuation of local resolvents,
generalized Gauss sums and anticyclotomic Hecke $L$-values of imaginary quadratic fields at inert
primes, preprint, arXiv:2507.08221.

\bibitem{DLMF} \emph{NIST Digital Library of Mathematical Functions}, Chapter~23,
\url{https://dlmf.nist.gov/23}.

\bibitem{GMR} P.~Guerzhoy, M.~H.~Mertens and L.~Rolen, Periodicities for Taylor coefficients of
half-integral weight modular forms, \emph{Pacific J.\ Math.} \textbf{307} (2020), 137--157.

\bibitem{HW} G.~H.~Hardy and E.~M.~Wright, \emph{An Introduction to the Theory of Numbers}, 6th
ed., Oxford University Press, 2008.

\bibitem{Hecke18} E.~Hecke, Eine neue Art von Zetafunktionen und ihre Beziehungen zur Verteilung der
Primzahlen, \emph{Math.\ Z.} \textbf{1} (1918), 357--376; \textbf{6} (1920), 11--51.

\bibitem{Hecke27} E.~Hecke, Theorie der Eisensteinschen Reihen h\"oherer Stufe und ihre Anwendung auf
Funktionentheorie und Arithmetik, \emph{Abh.\ Math.\ Sem.\ Univ.\ Hamburg} \textbf{5} (1927),
199--224.

\bibitem{Hurwitz} A.~Hurwitz, \"Uber die Entwicklungskoeffizienten der lemniskatischen Funktionen,
\emph{Math.\ Ann.} \textbf{51} (1899), 196--226.

\bibitem{IR} K.~Ireland and M.~Rosen, \emph{A Classical Introduction to Modern Number Theory},
2nd ed., Graduate Texts in Mathematics \textbf{84}, Springer, 1990.

\bibitem{Katz} N.~M.~Katz, Divisibilities, congruences, and Cartier duality, \emph{J.\ Fac.\ Sci.\
Univ.\ Tokyo Sect.\ IA Math.} \textbf{28} (1981), no.~3, 667--678.

\bibitem{KimLee} J.~Kim and Y.~Lee, $p$-adic properties for Taylor coefficients of half-integral
weight modular forms on $\Gamma_1(4)$, \emph{Taiwanese J.\ Math.} \textbf{27} (2023), no.~1,
23--38.

\bibitem{KM} C.~Krattenthaler and T.~W.~M\"uller, The congruence properties of Romik's sequence of
Taylor coefficients of Jacobi's theta function $\theta_3$, preprint, arXiv:2304.11471v4, 2024.

\bibitem{LarsonSmith} H.~Larson and G.~Smith, Congruence properties of Taylor coefficients of
modular forms, \emph{Int.\ J.\ Number Theory} \textbf{10} (2014), 1501--1518.

\bibitem{Mathlib} The mathlib Community, The Lean mathematical library, in: \emph{Proceedings of the
9th ACM SIGPLAN International Conference on Certified Programs and Proofs (CPP 2020)}, ACM, 2020,
367--381.

\bibitem{MT} B.~Mazur and J.~Tate, The $p$-adic sigma function, \emph{Duke Math.\ J.} \textbf{62}
(1991), no.~3, 663--688.

\bibitem{Lean4} L.~de~Moura and S.~Ullrich, The Lean~4 theorem prover and programming language, in:
\emph{Automated Deduction -- CADE 28}, Lecture Notes in Comput.\ Sci.\ \textbf{12699}, Springer,
2021, 625--635.

\bibitem{Nomoto} K.~Nomoto, The rank of a CM elliptic curve and a recurrence formula, preprint,
arXiv:2103.08947.

\bibitem{Romik} D.~Romik, The Taylor coefficients of the Jacobi theta constant $\theta_3$,
\emph{Ramanujan J.} \textbf{52} (2020), 275--290.

\bibitem{Scherer} R.~Scherer, Congruences modulo primes of the Romik sequence related to the Taylor
expansion of the Jacobi theta constant $\theta_3$, \emph{Ramanujan J.} \textbf{54} (2021),
427--448.

\bibitem{Silverman} J.~H.~Silverman, \emph{The Arithmetic of Elliptic Curves}, 2nd ed., Graduate
Texts in Mathematics \textbf{106}, Springer, 2009.

\bibitem{VZ} F.~Rodriguez Villegas and D.~Zagier, Square roots of central values of Hecke
$L$-series, in: \emph{Advances in Number Theory} (Kingston, ON, 1991), Oxford Univ.\ Press, 1993,
81--99.

\bibitem{Wakhare} T.~Wakhare, Romik's conjecture for the Jacobi theta function, \emph{J.\ Number
Theory} \textbf{215} (2020), 275--296.

\bibitem{Zagier} D.~Zagier, Elliptic modular forms and their applications, in: \emph{The 1-2-3 of
Modular Forms}, Universitext, Springer, 2008.

\end{thebibliography}

\end{document}
